\documentclass{article}
\usepackage{graphicx} % Required for inserting images
\usepackage{geometry}
\geometry{
a4paper,
total={170mm, 257mm},
left=20mm,
top=20mm,
}

\title{Labwork 1: Gradient Descent}

\begin{document}

\maketitle

\section{Algorithm implementation}

In this lab, we implement gradient descent from scratch. Here is the update rule:
$$x_{new} = x_{old} - r \cdot f'(x_{old})$$

where:
\begin{itemize}
\item $r$ is the learning rate
\item $f'(x)$ is the derivative of the function
\end{itemize}

The algorithm starts from an initial value and iteratively updates $x$.


\section{Effect of learning rate}

The learning rate $r$ plays a crucial role in the behavior of gradient descent:

\begin{itemize}
\item \textbf{Small learning rate ($r = 0.01$):}
Convergence is slow and requires many iterations.


\item \textbf{Moderate learning rate ($r = 0.1$):} 
Efficient convergence to the minimum.

\item \textbf{Large learning rate ($r = 1.0$):} 
The algorithm may miss the minimum or even diverge.


\end{itemize}

Thus, choosing an appropriate learning rate is important for both speed and stability.


\end{document}
